\tzpanchor{0369}
\tzpentryheader{0369}{Helicity Conservation Law}

\begingroup\small
\begingroup\scriptsize
\setlength{\tabcolsep}{4pt}
\renewcommand{\arraystretch}{0.92}
\noindent\begin{tabularx}{\linewidth}{@{}p{0.24\linewidth}p{0.36\linewidth}X@{}}
\textbf{UUID} & \multicolumn{2}{l@{}}{\tzpcode{9c01e38c-cfaf-58f7-9131-476cc2f25831}}\\
\textbf{PiID} & \tzpcode{5488204665213} & \textbf{TZPi}\quad \tzpcode{4}\quad \textbf{PiPE}\quad \tzpcode{376}\\
\textbf{RTEID} & \textbf{Poloidal $\theta$}\quad \tzpcode{1432.9089 rad} & \textbf{Toroidal $\phi$}\quad \tzpcode{02.3610 rad}\\
\textbf{TZPID} & \tzpcode{ID0369} & \textbf{Node Dependencies}\quad \tzpidref{0368}\\
\textbf{Registry} & \tzpcode{registry\_id\_0369} & \textbf{Scale}\quad diagnostic\\
\textbf{LOC Classification} & QA273 dimensionless analysis & QC statistical diagnostics\\
\textbf{arXiv Classification} & Primary: math-ph & Cross-list: nlin.PS, physics.data-an\\
\textbf{Primary Support} & Diagnostic Definition & Scaling Relation\\
\textbf{Closure Support} & Normalized Regime & \\
\end{tabularx}\par
\endgroup
\endgroup

\tzpsection{Canonical Statement}
Under ideal magneto-hydrodynamic evolution with KK coupling, the extended helicity defined above is conserved over time.

\tzpsection{Canonical Equation}
\[
\mathcal{H}=\int_{\Omega}\mathbf{A}\cdot\mathbf{B}\,d^{3}x
\]
\[
\frac{d\mathcal{H}}{dt}=0
\]

\begin{minipage}[t]{0.49\linewidth}
\tzpsection{Canonical Symbol Key}
\begingroup
\small
\raggedright
\textbf{Source equation symbols.}\par
\begin{itemize}[leftmargin=1.35em,itemsep=1pt,topsep=1pt,parsep=0pt]
    \item $\mathcal{H}$: source equation symbol.
    \item $\Omega$: domain, angular, or boundary operator associated with the equation.
    \item $A$: source equation symbol.
    \item $B$: magnetic field magnitude.
    \item $\Lambda$: cosmological or lemniscatic branch parameter in the entry relation.
    \item $\rho$: density or mass-density term used by the entry equation.
    \item $\mu$: source equation symbol.
\end{itemize}
\endgroup
\end{minipage}\hfill
\begin{minipage}[t]{0.47\linewidth}
\tzpsection{Wolfram Form}
\begingroup
\scriptsize
\raggedright
\textbf{Source Wolfram model.}\par
\begin{Verbatim}[fontsize=\tiny,breaklines=true,breakanywhere=true]
(* TZPID ID0369 generated source Wolfram model *)
ClearAll[K0369, Cr0369, T, R, A, W, I, N];
eq0369$1 = HoldForm[H == IntegralOmegaA*B*d^3*x];
eq0369$2 = HoldForm[dH/dt == 0];
eq0369$3 = HoldForm[Lambda == B^2/rhomuOmega^2];

K0369 = HoldForm[{eq0369$1, eq0369$2, eq0369$3}];

N = "normalized TRAWIN closure object for Helicity Conservation Law";
I = "Normalized Regime integration/comparison layer";
W = "Scaling Relation wave or operator carrier";
A = "Diagnostic Definition admissible action/amplitude condition";
R = "Scaling Relation rotation/curl preservation layer";
T = "Diagnostic Definition local source condition";

Cr0369[expr_] := HoldForm[N[I[W[A[R[T[expr]]]]]]];

Result0369 = Cr0369[K0369];
\end{Verbatim}
\endgroup
\end{minipage}

\par\vspace{0.45\baselineskip}
\noindent\begin{minipage}[t]{\linewidth}
\begingroup
\small
\raggedright
\textbf{TRAWIN Operator Key.}\par
\noindent\textbf{Time/localization operator:} $\mathsf{T}\!\left[\mathrm{local}\right]$ Diagnostic Definition local source condition.\par
\noindent\textbf{Rotation/curl operator:} $\mathsf{R}\!\left[\nabla\times\right]$ Scaling Relation rotation/curl preservation layer.\par
\noindent\textbf{Action/amplitude operator:} $\mathsf{A}\!\left[\mathrm{admissible}\right]$ Diagnostic Definition admissible action/amplitude condition.\par
\noindent\textbf{Wave/d'Alembertian operator:} $\mathsf{W}\!\left[\Box\right]$ Scaling Relation wave or operator carrier.\par
\noindent\textbf{Integration operator:} $\mathsf{I}\!\left[\int\right]$ Normalized Regime integration/comparison layer.\par
\noindent\textbf{Normalization operator:} $\mathsf{N}\!\left[\mathrm{normalized}\right]$ normalized TRAWIN closure object for Helicity Conservation Law.\par
\endgroup
\end{minipage}

\clearpage
\noindent\textbf{[ID 0369] Helicity Conservation Law}\par
\vspace{0.35em}
\tzpsection{Derivation Layer}
\paragraph{Primary Equation.}
\begin{equation}
m_{\mathrm{Goldstone}} = 0.
\end{equation}

\paragraph{Primary Derivation.}
Continuous symmetry breaking generates massless excitations.

Thus:
\begin{equation}
\boxed{
m_{\mathrm{Goldstone}} = 0.
}
\end{equation}

\paragraph{Interpretation.}
Broken symmetry produces gapless modes.

\par\vspace{0.35em}
\noindent\textbf{Derivable Consequences.}\quad
\begingroup
\footnotesize
\raggedright
The canonical equation anchors Helicity Conservation Law by connecting the source condition to the derivation layer and proof certificate.
\par\endgroup

\tzpsection{Equation Support}
\noindent\textbf{1. Diagnostic Definition}\par
\[
\mathcal{H}=\int_{\Omega}\mathbf{A}\cdot\mathbf{B}\,d^{3}x
\]
\noindent This defines the scalar diagnostic or conversion quantity named by the title.\par
\noindent\textbf{2. Scaling Relation}\par
\[
\frac{d\mathcal{H}}{dt}=0
\]
\noindent This elevates the quantity into a comparative scaling or correlation law.\par
\noindent\textbf{3. Normalized Regime}\par
\[
\Lambda=\frac{B^{2}}{\rho\mu\Omega^{2}}
\]
\noindent This closes the entry by normalizing the diagnostic against a stated reference regime.\par

\clearpage
\noindent\textbf{[ID 0369] Helicity Conservation Law}\par
\vspace{0.35em}
\tzpsection{Dictionary}
\begingroup\small
\tzpsection{Definition}
Helicity Conservation Law is a registry definition in Field Theory and Action Principles with secondary pressure from Manifold and Metric Dynamics.

% BEGIN TZP CONTEXTUAL MEANING
\tzpsection{Contextual Meaning}
\begingroup\footnotesize\setlength{\emergencystretch}{3em}\sloppy
\textbf{Formal statement:} Under ideal magneto-hydrodynamic evolution with KK coupling, the extended helicity defined above is conserved over time. \textbf{Contextual meaning:} The entry reads Helicity Conservation Law as a controlled technical headword whose equation layer, proof route, and registry role are interpreted together. \textbf{Derivable consequences:} The entry turns a named scalar quantity into a portable diagnostic that can anchor nearby comparisons. \textbf{Lemma support:} Diagnostic Definition; Scaling Relation; Normalized Regime.\par
\endgroup
% END TZP CONTEXTUAL MEANING

\tzpsection{Technical Interpretation}
Broken symmetry produces gapless modes.

\tzpsection{Equation Grounding}
Equation anchors: G(r, theta, phi) = g(theta, phi) (ratio r R ) to the power beta cos(ln ratio r R + phi sub 0); Lambda=B2rho Omega2Lambda; and frac B 2 rho,mu,Omega2 Lambda=rho. Proof route: show that the named quantity is dimensionless or consistently scaled in the regime stated by the entry. Formalization route: prove that the normalized diagnostic remains invariant under the intended rescaling.

\tzpsection{Interpretive Role}
The entry turns a named scalar quantity into a portable diagnostic that can anchor nearby comparisons.
\endgroup

\tzpsection{TZP Thesaurus}
\begin{enumerate}[leftmargin=1.6em,itemsep=6pt,topsep=4pt]
\item \textbf{Invulnerable Magnetic Topology}: The strict dynamic rule that, in an ideal fluid, the total number of magnetic knots cannot be created or destroyed.
\item \textbf{Constant Action Parameterization}: The normalization of the helicity integral to show that its value remains perfectly flat and unchanging over time.
\item \textbf{Dissipation-Free Flux Winding}: The condition where the resistance of the plasma is zero, preventing the magnetic field lines from cutting through one another.
\end{enumerate}

\tzpsection{Encyclopedia Mathematica}
\begingroup\footnotesize\sloppy
The visual plate below is reserved for the source-specific equation and progression graphic for [ID 0369]. It is included only as a companion to the canonical equation, not as a substitute for the formal text.
\par\endgroup
\ifTZPDarkEdition
\tzpdictionaryvisualband{D:/TZPID/kdp_workflow/build/tikz_figures/ID0369_dictionary_figure_dark.pdf}
\fi

\clearpage
\begingroup\tzpsupportsetup
\noindent\textbf{\fontsize{10.8pt}{12.8pt}\selectfont [ID 0369] Constructive Logic and Proof Certificate}\par
\noindent\textbf{\fontsize{10pt}{12pt}\selectfont Helicity Conservation Law}\par
\vspace{0.45em}
\begingroup
\footnotesize
\setlength{\parskip}{0.22em}
\setlength{\parindent}{0pt}
\noindent\textbf{Curry--Howard--Lambek Interpretation}\par
Under the Curry--Howard--Lambek correspondence, this registry entry is read as a typed proof
object: the stated source condition supplies the proposition, the derivation supplies the
morphism, and the proof certificate records the machine handles needed to check the obligation
in the formal registry.

\vspace{0.45em}
\noindent\textbf{CHL Proof}\par
\textbf{Statement:} under ideal magneto-hydrodynamic evolution with KK coupling, the extended helicity
defined above is conserved over time

\textbf{Logic Validation:}\\
\textbf{Proposition:} ``Helicity Conservation Law is valid when the total tensorial quantity is conserved under
the stated constraint.''\\
\textbf{Proof:} The source statement identifies a conserved balance law. Reading 0369 as a proposition,
validity follows by requiring the full stress, flux, or current contribution to be
included before the divergence or balance condition is evaluated.

\textbf{VERDICT:} \ensuremath{\checkmark}\ \textbf{TRUE} --- conserved total-structure validation

\textbf{Formal Status:} \ensuremath{\checkmark}\ THEORETICAL -- source-backed CHL proof obligation

\textbf{Physical Status:} THEORETICAL -- requires model-specific empirical interpretation
\par\vspace{0.45em}
\noindent{\tiny\textbf{Isabelle/HOL.} \tzpplaincode{physics\_validity\_from\_model, external\_validity\_package, registry\_number\_matches, equation\_count\_matches}}\par
\ifTZPDarkEdition
\tzpproofvisualband{D:/TZPID/kdp_workflow/build/tikz_figures/ID0369_proof_certificate_figure_dark.pdf}
\fi
\endgroup
\endgroup
